\begingroup%
\definecolor{ksBlue}{HTML}{3478B8}%
\definecolor{ksTeal}{HTML}{27877F}%
\definecolor{ksAmber}{HTML}{C58B28}%
\definecolor{ksInk}{HTML}{243447}%
\definecolor{ksSlate}{HTML}{667C92}%
\begin{tikzpicture}[
  x=1mm, y=1mm,
  card/.style={draw, fill=white, rounded corners=2mm, line width=.7pt,
    text=ksInk, align=flush center, font=\fontsize{9.5}{12.5}\selectfont,
    inner xsep=2mm, inner ysep=1.7mm, outer sep=.35pt},
  chapter/.style={card, minimum width=112mm},
  connector/.style={draw=ksInk, line width=.7pt, rounded corners=2mm},
  flow/.style={connector, -{Stealth[length=2.25mm,width=1.65mm]}}
]
  \node[card, draw=ksSlate, minimum width=86mm] (root) at (0,0)
    {问题根源\\
    Agent 缺少设计敏感度，也缺少设计载体\\
    把 Design 问题当作 Code 任务处理};
  \node[card, draw=ksBlue, line width=.9pt, minimum width=86mm] (center) at (0,-21)
    {\textbf{方案：Kernel Schema}\\
    跨硬件、跨语言的算子设计压缩表示};

  \node[chapter, draw=ksTeal] (w1) at (6,-41)
    {第五章 · 挑战一：数据与表示（已完成）\\
    算子压缩表示与知识蒸馏\\
    schema\_docs / template / extract / wiki};
  \node[chapter, draw=ksAmber] (w2) at (6,-61)
    {第六章 · 挑战二：生成与迭代（进行中）\\
    Planner–Coder 解耦闭环 · Step / Dispatch\\
    分离设计搜索与工程实现};
  \node[chapter, draw=ksSlate] (w3) at (6,-81)
    {第七章 · 挑战三：交付与维护（规划中）\\
    灰盒治理 · 分级异步验证 · 先验保序建模\\
    结合 Data 驱动与 Bench 验证};

  \coordinate (rail-top) at (-60,-21);
  \coordinate (rail-end) at (rail-top |- w3.west);
  \draw[flow] (root) -- (center);
  \draw[flow] (center.west) -- (rail-top) -- (rail-end) -- (w3.west);
  \draw[flow] (rail-top |- w1.west) -- (w1.west);
  \draw[flow] (rail-top |- w2.west) -- (w2.west);
\end{tikzpicture}%
\endgroup%
